\section{Module \texttt{button.c}}

Les boutons sont les déclencheurs primaires d'actions. Ce module centralise la logique de création des boutons, qu'ils soient textuels, iconiques ou un mélange des deux. L'apport majeur de ce wrapper est d'intégrer directement le raccordement des signaux (callbacks) lors de la création, éliminant ainsi le risque d'oublier l'appel à \texttt{g\_signal\_connect} plus loin dans le code.

\subsection{Fonction \texttt{create\_button}}

La fonction \texttt{create\_button} est conçue comme un point de rassemblement central pour toutes les formes de déclencheurs interactifs dans GTK. Historiquement, un développeur devait maîtriser différentes fonctions (\texttt{gtk\_button\_new\_with\_label}, \texttt{gtk\_button\_new\_with\_mnemonic}) et jongler avec des conteneurs supplémentaires s'il souhaitait ajouter une icône. Notre approche rationalise drastiquement ce processus.

Dans le corps de la fonction, la première étape logique consiste à discriminer le type de texte demandé. L'utilisation du drapeau booléen \texttt{use\_underline} oriente dynamiquement l'appel vers la fonction de création appropriée, permettant ainsi de supporter nativement les raccourcis clavier (mnémoniques) cruciaux pour l'accessibilité. 

La véritable complexité encapsulée par cette fonction apparaît lors de la gestion des icônes conjointes au texte. Si le champ \texttt{icon\_name} est renseigné, le code bascule d'une simple création textuelle vers la construction transparente d'une hiérarchie de widgets. Il instancie silencieusement une \texttt{GtkBox} interne, y insère une \texttt{GtkImage} configurée avec la taille précise exigée par l'API système, ajoute un \texttt{GtkLabel}, et intègre le tout comme unique enfant du bouton via \texttt{gtk\_button\_set\_child}. 

Cette automatisation cache complètement les rouages de la composition d'interface à l'utilisateur final. Enfin, la fonction s'assure d'injecter immédiatement les données métier (\texttt{user\_data}) et la fonction de rappel (\texttt{on\_clicked}) au signal d'activation, bouclant ainsi la création visuelle et fonctionnelle en un seul appel unifié.

\begin{lstlisting}[language=C]
GtkWidget *create_button(const button_config *config) {
    GtkWidget *button;
    const char *label;

    if (config == NULL) {
        return gtk_button_new();
    }

    label = config->label ? config->label : "";
    button = config->use_underline ? gtk_button_new_with_mnemonic(label) : gtk_button_new_with_label(label);
    gtk_button_set_has_frame(GTK_BUTTON(button), config->has_frame);
    gtk_button_set_can_shrink(GTK_BUTTON(button), config->can_shrink);

    if (config->icon_name && config->icon_name[0] != '\0') {
        GtkWidget *content = gtk_box_new(GTK_ORIENTATION_HORIZONTAL, 6);
        GtkWidget *icon = gtk_image_new_from_icon_name(config->icon_name);
        GtkWidget *text = config->use_underline ? gtk_label_new_with_mnemonic(label) : gtk_label_new(label);
        if (config->icon_size > 0) {
            gtk_image_set_pixel_size(GTK_IMAGE(icon), config->icon_size);
        }
        gtk_box_append(GTK_BOX(content), icon);
        gtk_box_append(GTK_BOX(content), text);
        gtk_button_set_child(GTK_BUTTON(button), content);
    }

    apply_theme_variant(button, config->theme_variant);
    apply_widget_style(button, &config->style);

    if (config->on_clicked) {
        g_signal_connect(button, "clicked", G_CALLBACK(config->on_clicked), config->user_data);
    }

    return button;
}
\end{lstlisting}

\paragraph{Macros associées :}
Les macros de variantes de thèmes (\texttt{BTN\_PRIMARY}, etc.) permettent de s'abstraire des chaînes de caractères brutes lors de la configuration. Bien que simples, elles jouent un rôle crucial pour la fiabilité du code : le compilateur signalera immédiatement une faute de frappe sur le nom de la macro, là où une erreur dans une chaîne littérale comme \texttt{"primay"} passerait inaperçue jusqu'à l'exécution, provoquant un rendu silencieusement incorrect de l'interface. Elles garantissent ainsi l'utilisation de noms de classes CSS cohérents avec la logique interne du module.

\begin{lstlisting}[language=C]
#define BTN_PRIMARY "primary"
#define BTN_SUCCESS "success"
#define BTN_DANGER  "danger"
#define BTN_WARNING "warning"
\end{lstlisting}

\subsection{Fonction \texttt{create\_radio\_button}}

Historiquement complexes à gérer, les boutons radio sont ici simplifiés par une interface déclarative directe. L'implémentation de \texttt{create\_radio\_button} reflète le nouveau standard de GTK4 qui a fusionné les anciens \texttt{GtkRadioButton} au profit d'un composant universel \texttt{GtkCheckButton} gérant des groupes logiques.

Le cœur de cette fonction repose sur un mécanisme d'enchaînement élégant. Au lieu de forcer le développeur à gérer lui-même une liste chaînée externe ou un objet groupe abstrait (comme le \texttt{GSList} sous GTK3), la configuration accepte simplement un pointeur vers un bouton précédemment créé (via le champ \texttt{group\_with}). La fonction vérifie la validité de ce pointeur et, si confirmé, utilise \texttt{gtk\_check\_button\_set\_group} pour raccorder mathématiquement le nouveau widget au groupe d'exclusion existant.

Ce paradigme réduit drastiquement les erreurs logiques de l'utilisateur. De plus, la définition immédiate de l'état par défaut (le booléen \texttt{is\_active}) et le branchement instantané de la callback sur l'événement \texttt{toggled} complètent l'initialisation. Le développeur n'a plus qu'à déclarer ses options l'une après l'autre, laissant le wrapper tisser les liens d'exclusivité en arrière-plan.

\begin{lstlisting}[language=C]
GtkWidget *create_radio_button(const radio_button_config *config) {
    GtkWidget *button = gtk_check_button_new_with_label(config && config->label ? config->label : "");

    if (config == NULL) {
        return button;
    }

    if (config->group_with != NULL) {
        gtk_check_button_set_group(GTK_CHECK_BUTTON(button), GTK_CHECK_BUTTON(config->group_with));
    }
    gtk_check_button_set_active(GTK_CHECK_BUTTON(button), config->is_active);
    apply_widget_style(button, &config->style);

    if (config->on_toggled) {
        g_signal_connect(button, "toggled", G_CALLBACK(config->on_toggled), config->user_data);
    }

    return button;
}
\end{lstlisting}

\subsection{APIs GTK utilisées}

\begin{itemize}
    \item \textbf{\texttt{gtk\_button\_new\_with\_label()}} : Crée un nouveau widget bouton contenant un label texte. Dans le module, elle est utilisée pour instancier le bouton de base lorsque l'option \texttt{use\_underline} est désactivée, offrant une création simplifiée de boutons textuels standards.
    \item \textbf{\texttt{gtk\_button\_new\_with\_mnemonic()}} : Crée un bouton dont le label contient un mnémonique (raccourci clavier) indiqué par un caractère souligné. Le wrapper y a recours lorsque \texttt{config->use\_underline} est vrai, facilitant ainsi l'accessibilité de l'interface au clavier.
    \item \textbf{\texttt{gtk\_button\_set\_child()}} : Définit le widget enfant unique du bouton. Elle permet au wrapper d'injecter une \texttt{GtkBox} personnalisée contenant à la fois une icône et un texte, dépassant ainsi la limitation des boutons à simple label.
    \item \textbf{\texttt{gtk\_image\_new\_from\_icon\_name()}} : Génère un widget d'image à partir du nom d'une icône présente dans le thème système. Elle est utilisée ici pour charger visuellement l'icône demandée dans la configuration avant son insertion dans le bouton.
    \item \textbf{\texttt{gtk\_check\_button\_set\_group()}} : Lie un bouton de type case à cocher à un groupe existant pour obtenir un comportement de sélection exclusive. Dans \texttt{create\_radio\_button}, cette API transforme un \texttt{GtkCheckButton} standard en un bouton radio fonctionnel.
    \item \textbf{\texttt{g\_signal\_connect()}} : Connecte une fonction de rappel à un signal spécifique d'un objet. Le module l'utilise pour brancher les signaux \texttt{clicked} (pour les boutons) et \texttt{toggled} (pour les radios) aux fonctions définies par le développeur applicatif.
\end{itemize}

\newpage
